% ====================================================================

% ====================================================================
\documentclass[article,paper=a4paper]{jlreq}

% jpreportパッケージの読み込み
\usepackage{jpreport}
\jpreportsetup{theoremwithin=none}
% ページレイアウトの設定
\usepackage{geometry}
\usepackage{enumitem}
\geometry{
    top=25mm,
    bottom=25mm,
    left=25mm,
    right=25mm
}
% TikZライブラリ
\usetikzlibrary{
    patterns,
    calc,
    arrows,
    decorations,
    angles,
    shapes.geometric,
    fit,
    knots
}
% TikZ用のスタイル
\tikzset{
  every path/.style={
    black,
    line width=1pt
  },
  every node/.style={
    transform shape,
    knot crossing,
    inner sep=1.5pt
  }
}
% ハイパーリンクの設定
\usepackage[unicode]{hyperref}
\usepackage{bookmark}
\usepackage[nameinlink]{cleveref}

\hypersetup{
    colorlinks=true,
    linkcolor=black,
    urlcolor=blue,
    citecolor=blue,
    pdftitle=CM環の幾何学的直感について,
    pdfauthor=森口湧成
}

% 数学演算子の定義


% cleveref用の日本語設定
\crefname{thm}{定理}{定理}
\crefname{lem}{補題}{補題}
\crefname{cor}{系}{系}
\crefname{ex}{例}{例}
\crefname{def}{定義}{定義}
\crefname{prop}{命題}{命題}
\crefname{prob}{問題}{問題}

% 色彩の定義
\definecolor{yukari}{RGB}{196, 163, 191}

% 文書情報
\title{CM環の幾何学的直感について}
\author{森口湧成\thanks{徳島大学 理工学部 理工学科 数理科学コース3年}}
\date{\today}

\begin{document}
\pagestyle{empty}

% --- タイトル ------------------------------------------------------
\maketitle
\thispagestyle{empty}
\begin{abstract}
    初めましての方は初めまして. そうでない方はこんにちは. 今回, 初めて puits に寄稿させていただきます第42回数理の翼夏季セミナー参加・第46回数理の翼夏季セミナーのスタッフの森口湧成と申します. 普段は代数学を中心に勉強しており, 最近は Condensed Mathematics を中心として数論幾何方面の勉強をしています. 数学系の話はTwitter(\href{https://x.com/akari0koutya}{@akari0koutya})で, 「上伊那ぼたん, 酔へる姿は百合の花」や「超かぐや姫！」など, 漫画やアニメの話, 湧源とかの話, 数学以外の話全ては 
Twitter(\href{https://x.com/Charle__F}{@Charle\_\_F})
 でしています. よろしければそちらも見ていただけると嬉しいです.

    \vspace{1\baselineskip}

    今回は, 工藤さんが湧源のB2,3あたりの数学徒が puits に寄稿しないことを嘆かれていたので書いてみました. 内容は, 代数学の中の Cohen-Macaulay 環 (CM環) というクラスの環の幾何学的な直感についてです. CM環は正則環よりも広いクラスの環でありながら, 非孤立素因子を持たないなどの整った振る舞いをすることが知られています. 今回の記事では, 具体例を通して CM環の幾何学的なイメージを掴むことを目指します. 尚, 雰囲気を掴むことを目的としているので $\operatorname{Ext}$ といった derived な道具は一切使わないことにします. そのため, 関連する scheme 論的な話である dualizing sheaf の話もしませんし, 証明はすべて省略します.

\end{abstract}
%% 本文の内容をここに記述
% ====================================================================
\begin{rem}[本稿で使う主な記号]
    本稿では $k$ を体とする. また, ネーター環 $A$, イデアル $I\subset A$, 有限生成
    $A$-\hspace{0pt}加群 $M$ に対して, 記号の意味を次のように約束する.
    \begin{enumerate}[label=(\arabic*)\,]
        \item $\Spec(A)$ は $A$ の素イデアル全体からなる空間である.
        \item $V(I)$ は $\Spec(A)$ の閉集合
        \[
            V(I):=\{\mathfrak{p}\in \Spec(A)\mid I\subset \mathfrak{p}\}
        \]
        を表す. たとえば $V(x)$ は $V((x))$ の略記である.
        \item $A_{\mathfrak{p}}, M_{\mathfrak{p}}$ は素イデアル $\mathfrak{p}$ による局所化である.
        特に $\mathfrak{m}$ が極大イデアルのとき, $A_{\mathfrak{m}}$ はその点における局所環である.
        \item $\Supp(M)$ は $M$ の台集合,
        \[
            \Supp(M):=\{\mathfrak{p}\in \Spec(A)\mid M_{\mathfrak{p}}\neq 0\}
        \]
        であり, 本稿では
        \[
            \dim(M):=\dim(\Supp(M))
        \]
        と約束する. 特に環 $A$ 自身についての $\dim(A)$ は通常の Krull 次元である.
        \item $\Ass_A(M)$ は $M$ の素因子全体である.
        \item $\sqrt{I}$ はイデアル $I$ の根基, $A_{\mathrm{red}}:=A/\sqrt{(0)}$ は
        $A$ の還元環である. つまり nilpotent 元をすべて 0 にした環である.
        \item $\operatorname{ht}(I)$ はイデアル $I$ の高さであり, $I$ を含む素イデアルの高さの最小値である.
        \item 極大イデアル $\mathfrak{m}$ に対して $\mathfrak{m}/\mathfrak{m}^2$ は
        cotangent space を表す. その双対が Zariski 接空間である.
    \end{enumerate}
\end{rem}

\begin{dfn}\label{def:depth}
    \quad $(A,\mathfrak{m})$ : ネーター局所環, $M \neq 0$ : 有限生成 $A$-\hspace{0pt}加群. \\
    $M$ の\textbf{深さ (depth)} とは, $\mathfrak{m}$ の元からなる $M$-正則列の長さの上限
    \[
        \depth_A(M)
        :=
        \sup\left\{\, r \in \mathbb{Z}_{\geq 0} \;\middle|\; \text{$\mathfrak{m}$ の元からなる長さ $r$ の $M$-正則列が存在する} \,\right\}
    \]
    のことである. 本稿では環が明らかなとき $\depth_A(M)$ を単に $\depth(M)$ と書く.
\end{dfn}

\begin{dfn}\label{def:cm}
    \quad $A$ : ネーター局所環, \quad $M$ : 有限生成 $A$-\hspace{0pt}加群.
    \begin{enumerate}[label=(\arabic*)\,]
        \item $M$ が \textbf{Cohen-Macaulay加群} であるとは, $\depth(M) = \dim(M)$ であることをいう.
        \item $A$ が \textbf{Cohen-Macaulay局所環} であるとは, $A$ を $A$-加群とみたときに Cohen-Macaulay加群 であることをいう.
        \item ネーター環 $A$ が \textbf{Cohen-Macaulay環} であるとは, 任意の極大イデアル $\mathfrak{m}$ に対して局所環 $A_{\mathfrak{m}}$ が Cohen-Macaulay局所環であることをいう.
    \end{enumerate}
    Cohen-Macaulay加群, Cohen-Macaulay局所環, Cohen-Macaulay環 をそれぞれ
    \textbf{CM加群}, \textbf{CM局所環}, \textbf{CM環} と略すことがある. ここでもそれを採用する.
\end{dfn}

最初の例を読むために, いくつか基本語を先に定義しておく.

\begin{dfn}\label{def:regular-sequence}
    \quad $A$ : ネーター環, $M$ : 有限生成 $A$-\hspace{0pt}加群, $x_1,\dots,x_r \in A$.\\
    列 $x_1,\dots,x_r$ が \textbf{$M$-正則列} であるとは, 各 $i=1,\dots,r$ に対して
    $x_i$ が
    \[
        M/(x_1,\dots,x_{i-1})M
    \]
    上の零因子でなく, さらに
    \[
        M/(x_1,\dots,x_r)M \neq 0
    \]
    であることをいう. 特に $M=A$ のとき \textbf{$A$-正則列} という.
\end{dfn}

\begin{dfn}\label{def:regular-local}
    \quad $(A,\mathfrak{m})$ : ネーター局所環. \\
    $A$ が \textbf{正則局所環} であるとは, 極大イデアル $\mathfrak{m}$ の最小生成元の個数が
    $\dim A$ に一致することをいう.
\end{dfn}

\begin{rem}
    正則局所環の幾何学的意味は「その点が滑らかである」ことである. 実際, 剰余体を
    $k=A/\mathfrak{m}$ とすると, 極大イデアル $\mathfrak{m}$ の最小生成元の個数は
    \[
        \dim_k(\mathfrak{m}/\mathfrak{m}^2)
    \]
    に等しい. これはその点のまわりの一次近似, すなわち cotangent space の次元であり,
    その双対が Zariski接空間になる. したがって正則局所環とは, 空間の次元と接空間の次元が
    一致している局所環, すなわち幾何学的には特異でない点に対応する局所環である.
\end{rem}

定義\ref{def:depth}, \ref{def:regular-sequence} によれば, $\depth(M)$ は
$M$ 上で零因子にならない元を何本続けて取れるかを測る不変量である. 一般には
\[
    \depth(M)\le \dim(M)
\]
が成り立つので, 定義\ref{def:cm} の CM条件 $\depth(M)=\dim(M)$ は
「次元の分だけ正則列を取れる」という意味で, 深さが可能な限り大きいことを表している.
環の場合には $M=A$ とみればよい. したがって CM環とは, すべての極大イデアルで局所化したときに
「深さが次元にぴったり一致する」ような環のことである.

CM性は正則性より弱い条件である. 正則局所環ならば必ず CM局所環だが, 逆は一般には成り立たない.
したがって CM環とは「滑らかである」ことではなく, 特異ではあっても depth と次元がよく揃った
整った振る舞いを許す概念だと理解するとよい.

\begin{rem}
    ここで正則局所環と CM局所環の違いをまとめておく. 正則局所環が表しているのは
    「接空間の次元がちょうど空間の次元に等しい」という滑らかさである. 他方, CM局所環が
    表しているのは
    \[
        \depth(A)=\dim(A)
    \]
    という深さの最大性であり, こちらは特異点を許す. したがって
    \[
        \text{正則局所環} \Longrightarrow \text{CM局所環}
    \]
    だが逆は一般には成り立たない. 正則局所環が「滑らかさ」を測る概念であるのに対して,
    CM局所環は「特異ではあっても余分な非孤立素因子を持たず, 正則列が最大限に取れる」
    という整い方を測る概念だと見ると分かりやすい.
\end{rem}

まず, その代表例として
\[
    R = k[x,y,z]/(xy)
\]
を考える. ここで $k[x,y,z]$ は体 $k$ 上の 3 変数多項式環であり,
$(xy)$ は元 $xy$ が生成する主イデアルである. この環は 2 枚の平面が交線に沿って交わる
特異な空間を与えるが, それでも CM環になる.

\begin{prop}\label{prop:cm-quotient}
    \quad $(A,\mathfrak{m})$ : CM局所環, $x_1,\dots,x_r \in \mathfrak{m}$ : $A$-正則列. \\
    このとき, $A/(x_1,\dots,x_r)$ も CM局所環である.
\end{prop}

\begin{cor}\label{cor:regular-local-cm}
    正則局所環は CM局所環である. したがって, 正則局所環を正則列で割って得られる局所環,
    特に 1 本の非零因子で割って得られる局所環は CM局所環である.
\end{cor}

まずは 命題\ref{prop:cm-quotient} と 系\ref{cor:regular-local-cm} を使って, $R$ が本当に CM環であることを確認する.
任意の極大イデアル $\mathfrak{m}\subset R$ に対し, その逆像を
$\widetilde{\mathfrak{m}}\subset k[x,y,z]$ とおくと
\[
    R_{\mathfrak{m}}\cong k[x,y,z]_{\widetilde{\mathfrak{m}}}/(xy)
\]
である. ここで $\widetilde{\mathfrak{m}}$ は自然な全射
\[
    k[x,y,z]\twoheadrightarrow R
\]
による $\mathfrak{m}$ の逆像であり, $k[x,y,z]_{\widetilde{\mathfrak{m}}}$ は
周囲の 3 次元アフィン空間の対応する点における局所環である.
また $k[x,y,z]_{\widetilde{\mathfrak{m}}}$ は定義\ref{def:regular-local} の意味で正則局所環であり,
$k[x,y,z]_{\widetilde{\mathfrak{m}}}$ は整域なので $xy$ は零因子ではない.
したがって $R_{\mathfrak{m}}$ は「正則局所環 $k[x,y,z]_{\widetilde{\mathfrak{m}}}$ を
1 本の非零因子 $xy$ で割って得られる局所環」とみなせる.
まず 系\ref{cor:regular-local-cm} により
$k[x,y,z]_{\widetilde{\mathfrak{m}}}$ は CM局所環であり, さらに
命題\ref{prop:cm-quotient} を長さ 1 の正則列 $(xy)$ に適用すると
各 $R_{\mathfrak{m}}$ も CM局所環になる. よって $R$ は CM環である.

幾何学的には, $R$ は図\ref{fig:cm-overview} のように 3 次元空間の中で
2 枚の平面 $V(x),V(y)$ が交線 $V(x,y)$ に沿って交わる空間である.

\vspace{2\baselineskip}

\begin{figure}[H]
    \centering
    \begin{tikzpicture}[
        line join=round,
        x={(1.18cm,0cm)},
        y={(-0.74cm,0.48cm)},
        z={(0cm,1.15cm)},
        info/.style={
            font=\small,
            fill=white,
            fill opacity=0.9,
            text opacity=1,
            rounded corners=2pt,
            inner sep=2pt
        },
        callout/.style={
            ->,
            line width=0.9pt
        }
    ]
        \def\xmin{-2.15}
        \def\xmax{2.15}
        \def\ymin{-1.85}
        \def\ymax{1.85}
        \def\zmax{3.05}
        \def\zs{1.45}

        % ambient affine 3-space
        \draw[gray!55, dashed] (\xmin,\ymin,0) -- (\xmax,\ymin,0) -- (\xmax,\ymax,0) -- (\xmin,\ymax,0) -- cycle;
        \draw[gray!55, dashed] (\xmin,\ymin,\zmax) -- (\xmax,\ymin,\zmax) -- (\xmax,\ymax,\zmax) -- (\xmin,\ymax,\zmax) -- cycle;
        \draw[gray!55, dashed] (\xmin,\ymin,0) -- (\xmin,\ymin,\zmax);
        \draw[gray!55, dashed] (\xmax,\ymin,0) -- (\xmax,\ymin,\zmax);
        \draw[gray!55, dashed] (\xmax,\ymax,0) -- (\xmax,\ymax,\zmax);
        \draw[gray!55, dashed] (\xmin,\ymax,0) -- (\xmin,\ymax,\zmax);

        % horizontal slice
        \path[fill=gray!12, draw=gray!55, fill opacity=0.45]
            (\xmin,\ymin,\zs) -- (\xmax,\ymin,\zs) -- (\xmax,\ymax,\zs) -- (\xmin,\ymax,\zs) -- cycle;

        % coordinate planes
        \path[fill=blue!18, draw=blue!65!black, fill opacity=0.55]
            (0,\ymin,0) -- (0,\ymax,0) -- (0,\ymax,\zmax) -- (0,\ymin,\zmax) -- cycle;
        \path[fill=green!18, draw=green!50!black, fill opacity=0.55]
            (\xmin,0,0) -- (\xmax,0,0) -- (\xmax,0,\zmax) -- (\xmin,0,\zmax) -- cycle;

        % intersection and slice traces
        \draw[red!70!black, line width=1.6pt] (0,0,0) -- (0,0,\zmax);
        \draw[blue!50!black, line width=1.2pt] (0,\ymin,\zmax) -- (0,\ymax,\zmax);
        \draw[blue!50!black, line width=1.2pt] (0,\ymin,0) -- (0,\ymin,\zmax);
        \draw[blue!50!black, line width=1.2pt] (0,\ymin,\zs) -- (0,\ymax,\zs);
        \draw[green!50!black, line width=1.2pt] (\xmin,0,0) -- (\xmin,0,\zmax);
        \draw[green!50!black, line width=1.2pt] (\xmin,0,\zs) -- (\xmax,0,\zs);
        \fill[black] (0,0,\zs) circle (2.2pt);

        \coordinate (vxTarget) at (0,1.28,2.45);
        \coordinate (vyTarget) at (1.82,0,2.52);
        \coordinate (lineTarget) at (0,0,2.70);
        \coordinate (pointTarget) at (0,0,\zs);
        \coordinate (sliceTarget) at (1.18,0.40,\zs);

        \node[info, blue!65!black] (vxLabel) at (-3.95,0.90,2.78) {$V(x)$};
        \draw[callout, blue!65!black] (vxLabel.east) -- (vxTarget);

        \node[info, green!50!black] (vyLabel) at (5.00,0.35,3.00) {$V(y)$};
        \draw[callout, green!50!black] (vyLabel.west) -- (vyTarget);

        \node[info, red!70!black] (lineLabel) at (-0.10,2.55,4.00) {$V(x,y)$};
        \draw[callout, red!70!black] (lineLabel.south) -- (lineTarget);

        \node[info, anchor=west] (pointLabel) at (5.35,-0.92,1.02) {交線上の点 $\mathfrak{m}_c$};
        \draw[callout] (pointLabel.west) -- (pointTarget);

        \node[info, align=left] (sliceLabel) at (5.55,-2.05,2.78) {断面 $z=c$ では\\$x=0$ と $y=0$ が交差};
        \draw[callout, black!70] (sliceLabel.west) -- (sliceTarget);
    \end{tikzpicture}
    \caption{$R=k[x,y,z]/(xy)$ の概形的形. 2 つの座標平面 $V(x),V(y)$ が交線 $V(x,y)$ に沿って交わっている.}
    \label{fig:cm-overview}
\end{figure}

この例は, 正則局所環と CM局所環の差を具体的にも見せている. 交線上の点
\[
    \mathfrak{m}_c=(x,y,z-c)\subset R
\]
は, $c\in k$ に対して交線 $V(x,y)$ 上の閉点 $(0,0,c)$ に対応する極大イデアルである.
添字 $c$ はその点の $z$ 座標を表している. この点で局所化すると
\[
    \dim(R_{\mathfrak{m}_c})=2
\]
であるが, 極大イデアル
$\mathfrak{m}_cR_{\mathfrak{m}_c}$ の一次部分は $x,y,z-c$ の 3 方向を持つ.
実際, 関係式 $xy=0$ は 2 次式であって 1 次の関係を与えないので,
\[
    \dim_{k}\bigl(\mathfrak{m}_c/\mathfrak{m}_c^2\bigr)=3
\]
となる. したがって $R_{\mathfrak{m}_c}$ は正則局所環ではない. 幾何学的には,
2 枚の平面が交差しているため, 交線上では接空間が 2 次元でなく 3 次元に膨らんでいるのである.
それでも $R_{\mathfrak{m}_c}$ は CM局所環である. ここに「正則ではないが CM である」という現象が現れている.

比較のため
\[
    A_{\mathrm{reg}}=k[u,v]_{(u,v)}
\]
を固定する. これは平面 $\mathbb{A}^2_k$ の原点における局所環であり, その極大イデアルは
\[
    \mathfrak{m}_{\mathrm{reg}}=(u,v)A_{\mathrm{reg}}
\]
である.
\[
    \dim A_{\mathrm{reg}}=2,\qquad \dim_k(\mathfrak{m}_{\mathrm{reg}}/\mathfrak{m}_{\mathrm{reg}}^2)=2
\]
だから正則局所環である. これに対して $R_{\mathfrak{m}_c}$ では
\[
    \dim R_{\mathfrak{m}_c}=2,\qquad \dim_k(\mathfrak{m}_c/\mathfrak{m}_c^2)=3
\]
となる. 図\ref{fig:regular-vs-cm} はこの差を模式的に表している.
\vspace{1\baselineskip}
\begin{figure}[H]
    \centering
    \begin{tikzpicture}[
        line join=round,
        x={(1.18cm,0cm)},
        y={(-0.74cm,0.48cm)},
        z={(0cm,1.15cm)},
        info/.style={
            font=\small,
            fill=white,
            fill opacity=0.9,
            text opacity=1,
            rounded corners=2pt,
            inner sep=2pt
        },
        callout/.style={
            ->,
            line width=0.9pt
        },
        vec/.style={
            ->,
            line width=1pt
        }
    ]
        \def\xmin{-1.75}
        \def\xmax{1.75}
        \def\ymin{-1.55}
        \def\ymax{1.55}
        \def\zmax{2.70}
        \def\zs{1.35}
        \def\ptz{\zs}

        % regular local ring
        \begin{scope}[shift={(-4.76cm,0cm)}]
            \draw[gray!55, dashed] (\xmin,\ymin,0) -- (\xmax,\ymin,0) -- (\xmax,\ymax,0) -- (\xmin,\ymax,0) -- cycle;
            \draw[gray!55, dashed] (\xmin,\ymin,\zmax) -- (\xmax,\ymin,\zmax) -- (\xmax,\ymax,\zmax) -- (\xmin,\ymax,\zmax) -- cycle;
            \draw[gray!55, dashed] (\xmin,\ymin,0) -- (\xmin,\ymin,\zmax);
            \draw[gray!55, dashed] (\xmax,\ymin,0) -- (\xmax,\ymin,\zmax);
            \draw[gray!55, dashed] (\xmax,\ymax,0) -- (\xmax,\ymax,\zmax);
            \draw[gray!55, dashed] (\xmin,\ymax,0) -- (\xmin,\ymax,\zmax);
            \path[fill=gray!12, draw=gray!55, fill opacity=0.45]
                (\xmin,\ymin,\zs) -- (\xmax,\ymin,\zs) -- (\xmax,\ymax,\zs) -- (\xmin,\ymax,\zs) -- cycle;
            \path[fill=blue!18, draw=blue!65!black, fill opacity=0.55]
                (0,\ymin,0) -- (0,\ymax,0) -- (0,\ymax,\zmax) -- (0,\ymin,\zmax) -- cycle;
            \draw[blue!50!black, line width=1.2pt] (0,\ymin,\zmax) -- (0,\ymax,\zmax);
            \draw[blue!50!black, line width=1.2pt] (0,\ymin,0) -- (0,\ymin,\zmax);
            \draw[blue!50!black, line width=1.2pt] (0,\ymin,\zs) -- (0,\ymax,\zs);
            \fill (0,0,\ptz) circle (2.2pt);
            \draw[vec, blue!65!black] (0,0,\ptz) -- (0,0.95,\ptz);
            \draw[vec, blue!65!black] (0,0,\ptz) -- (0,0,\ptz+0.95);

            \node[info] at (0,0,3.65) {$A_{\mathrm{reg}}=k[u,v]_{(u,v)}$};
            \node[info, align=center] at (0,0,-1.25)
                {滑らかな点\\$\dim=2,\ \dim_k(\mathfrak m/\mathfrak m^2)=2$};
            \node[font=\scriptsize] at (0.10,1.10,\ptz+0.1) {$u$};
            \node[font=\scriptsize, blue!65!black] at (0.10,0.08,\ptz+1.08) {$v$};
        \end{scope}

        % CM but non-regular local ring
        \begin{scope}[shift={(4.76cm,0cm)}]
            \draw[gray!55, dashed] (\xmin,\ymin,0) -- (\xmax,\ymin,0) -- (\xmax,\ymax,0) -- (\xmin,\ymax,0) -- cycle;
            \draw[gray!55, dashed] (\xmin,\ymin,\zmax) -- (\xmax,\ymin,\zmax) -- (\xmax,\ymax,\zmax) -- (\xmin,\ymax,\zmax) -- cycle;
            \draw[gray!55, dashed] (\xmin,\ymin,0) -- (\xmin,\ymin,\zmax);
            \draw[gray!55, dashed] (\xmax,\ymin,0) -- (\xmax,\ymin,\zmax);
            \draw[gray!55, dashed] (\xmax,\ymax,0) -- (\xmax,\ymax,\zmax);
            \draw[gray!55, dashed] (\xmin,\ymax,0) -- (\xmin,\ymax,\zmax);
            \path[fill=gray!12, draw=gray!55, fill opacity=0.45]
                (\xmin,\ymin,\zs) -- (\xmax,\ymin,\zs) -- (\xmax,\ymax,\zs) -- (\xmin,\ymax,\zs) -- cycle;
            \path[fill=blue!18, draw=blue!65!black, fill opacity=0.55]
                (0,\ymin,0) -- (0,\ymax,0) -- (0,\ymax,\zmax) -- (0,\ymin,\zmax) -- cycle;
            \path[fill=green!18, draw=green!50!black, fill opacity=0.55]
                (\xmin,0,0) -- (\xmax,0,0) -- (\xmax,0,\zmax) -- (\xmin,0,\zmax) -- cycle;

            \draw[red!70!black, line width=1.6pt] (0,0,0) -- (0,0,\zmax);
            \draw[blue!50!black, line width=1.2pt] (0,\ymin,\zmax) -- (0,\ymax,\zmax);
            \draw[blue!50!black, line width=1.2pt] (0,\ymin,0) -- (0,\ymin,\zmax);
            \draw[blue!50!black, line width=1.2pt] (0,\ymin,\zs) -- (0,\ymax,\zs);
            \draw[green!50!black, line width=1.2pt] (\xmin,0,0) -- (\xmin,0,\zmax);
            \draw[green!50!black, line width=1.2pt] (\xmin,0,\zs) -- (\xmax,0,\zs);
            \fill[black] (0,0,\ptz) circle (2.2pt);
            \draw[vec] (0,0,\ptz) -- (1.10,0,\ptz);
            \draw[vec] (0,0,\ptz) -- (0,0.95,\ptz);
            \draw[vec] (0,0,\ptz) -- (0,0,\ptz+0.95);

            \node[info] at (0,0,3.65) {$R_{\mathfrak{m}_c}$};
            \node[info, align=center] at (0,0,-1.25)
                {CM だが正則でない\\$\dim=2,\ \dim_k(\mathfrak m_c/\mathfrak m_c^2)=3$};
            \node[font=\scriptsize] at (1.27,0.10,\ptz+0.1) {$x$};
            \node[font=\scriptsize] at (0.10,1.10,\ptz+0.1) {$y$};
            \node[font=\scriptsize] at (-0.28,0.08,\ptz+1.08) {$z-c$};
        \end{scope}
    \end{tikzpicture}
    \vspace{2\baselineskip}
    \caption{左は正則局所環 $A_{\mathrm{reg}}$ の模式図で, 切断面は 1 本の直線に見え, 接方向は 2 本しかない. 右は $R_{\mathfrak{m}_c}$ の模式図で, 切断面が十字に見え, 接方向が 3 本に増えている. これが「CM だが正則でない」という差である.}
    \label{fig:regular-vs-cm}
\end{figure}
左の正則局所環では接空間の次元が空間の次元と一致しているが, 右の $R_{\mathfrak{m}_c}$ では
2 枚の平面が交差するために余分な接方向が現れ, そこが正則性の破れとして見えている.

次に 定義\ref{def:cm} の等式 $\depth=\dim$ がこの例でどう見えるかを書く.
実際, $\dim(R)=2$ である一方, $z$ は $R$ の零因子ではなく,
\[
    R/(z)\cong k[x,y]/(xy)
\]
において $x+y$ も零因子ではない. したがって $z,x+y$ は $R$-正則列になり,
\[
    \depth(R)=2=\dim(R)
\]
となる. これがこの環が CM環であることの, depth を通した具体的な見え方である.

この計算は, CM局所環の一般論ともよく整合している. 局所環では, CM性は
「次元を測る元たちが正則列になること」と言い換えられる.

\begin{dfn}\label{def:parameter-system}
    \quad $(A,\mathfrak{m})$ : ネーター局所環, $\dim A = d$.\\
    元 $x_1,\dots,x_d \in \mathfrak{m}$ が \textbf{パラメータ系} であるとは,
    イデアル $(x_1,\dots,x_d)$ が $\mathfrak{m}$-\hspace{0pt}準素であることをいう.
\end{dfn}

\begin{prop}\label{prop:cm-regular-sequence}
    \quad $(A,\mathfrak{m})$ : ネーター局所環. \quad このとき次は同値である.
    \begin{enumerate}[label=(\arabic*)\,]
        \item $A$ は CM局所環である.
        \item あるパラメータ系が $A$-正則列である.
        \item 任意のパラメータ系が $A$-正則列である.
    \end{enumerate}
\end{prop}

言い換えると, CM局所環では「次元を削るための元」が零因子を作らずに並ぶ.
上で見た正則列 $z,x+y$ は, その感覚をこの具体例で見せている.

この環の幾何は非常に見やすい. 実際,
\[
    (xy)=(x)\cap(y)
\]
が成り立つので, $R$ に対応する空間は $k$ 上の 3 次元アフィン空間の中の 2 つの座標平面
$V(x)$ と $V(y)$ の合併になる. 言い換えると, 方程式 $xy=0$ は
「$x=0$ または $y=0$」を表しており, 空間全体は 2 枚の平面が交線 $V(x,y)$ に沿って
貼り合わさった形をしている.

では, 交線の generic point に対応する $(x,y)$ がなぜ素因子にならないのだろうか.
ここで使う一般事実は次である.

\begin{prop}\label{prop:cm-associated-primes}
    \quad $A$ : ネーター局所環, $M$ : CM加群. \\
    このとき, $M$ の任意の素因子 $\mathfrak{p} \in \Ass_{A}(M)$ に対して次が成り立つ:
    \[\dim(A/\mathfrak{p}) = \dim(M) = \depth(M).\]
\end{prop}

\begin{rem}
    上の命題から, 局所環上の CM加群 $M$ は非孤立素因子を持たないことがわかる.
\end{rem}

\begin{dfn}\label{def:unmixedness}
    \quad $A$ : ネーター環. \\
    $A$ が \textbf{純正定理}を満たすとは, 任意の整数 $r\ge 0$ と
    任意の $r$ 個の元 $a_1,\dots,a_r\in A$ に対し,
    \[
        I=(a_1,\dots,a_r),\qquad \operatorname{ht}(I)=r
    \]
    ならば, $A/I$ の任意の素因子 $\mathfrak{p}\in \Ass_A(A/I)$ が
    \[
        \operatorname{ht}(\mathfrak{p})=r
    \]
    を満たすことである. 言い換えると, 期待される余次元 $r$ で切った閉部分
    $V(I)$ には余分な非孤立素因子が現れず, すべての素因子が同じ余次元 $r$ をもつ. 
\end{dfn}

\begin{thm}\label{thm:cm-iff-unmixed}
    \quad $A$ : ネーター環. \quad このとき次は同値である.
    \begin{enumerate}
        \item $A$ は CM環である. 
        \item $A$ は 定義\ref{def:unmixedness} の意味で純正定理を満たす. 
    \end{enumerate}
\end{thm}

\begin{cor}\label{cor:purity-local}
    \quad $(A,\mathfrak{m})$ : CM局所環, $x_1,\dots,x_r$ : $A$-正則列,
    $I=(x_1,\dots,x_r)$ . \\
    このとき $A/I$ は CM局所環であり, そのすべての素因子は同じ次元
    \[
        \dim(A)-r
    \]
    をもつ. 特に $A/I$ は非孤立素因子を持たない. 
\end{cor}

\begin{rem}
    これは「正則列で切ってできる閉部分は, 期待される次元を保ち, 余分な非孤立素因子を持たない」
    という純性現象である. 有限型の幾何を思い浮かべれば, 各既約成分は一様な余次元 $r$ をもち,
    切ったことで突然より高い余次元の成分が湧いてくることはない. ここに
    定理\ref{thm:cm-iff-unmixed} の局所的な現れがある.
\end{rem}

ここで重要なのは, 交線 $V(x,y)$ は 2 つの既約成分の交わりではあるが,
それ自身が新しい既約成分ではないという点である. したがって素因子は
各既約成分の generic point に対応する $(x)$ と $(y)$ だけであり, 交線の generic point に対応する
$(x,y)$ が素因子になるわけではない. この点で 命題\ref{prop:cm-associated-primes} が働く.
実際, もし $(x,y)$ が $R$ の素因子だとすると, $(x,y)$ を含む極大イデアル $\mathfrak{m}$ を取れば
\[
    (x,y)R_{\mathfrak{m}} \in \Ass_{R_{\mathfrak{m}}}(R_{\mathfrak{m}})
\]
となる. ところが $R_{\mathfrak{m}}$ は CM局所環なので, 命題\ref{prop:cm-associated-primes} により
任意の素因子 $\mathfrak{p}\in \Ass_{R_{\mathfrak{m}}}(R_{\mathfrak{m}})$ は
\[
    \dim(R_{\mathfrak{m}}/\mathfrak{p})=\dim(R_{\mathfrak{m}})=2
\]
を満たさなければならない. しかし
\[
    R_{\mathfrak{m}}/(x,y)R_{\mathfrak{m}}
    \cong (R/(x,y))_{\mathfrak{m}/(x,y)}
    \cong k[z]_{\mathfrak{m}/(x,y)}
\]
であるから
\[
    \dim\bigl(R_{\mathfrak{m}}/(x,y)R_{\mathfrak{m}}\bigr)=1
\]
となって矛盾する. したがって $(x,y)$ は $R$ の素因子ではない.
CM環でもこのように, 各極大イデアルで局所化した CM局所環の性質を通して見ると,
成分どうしが交わっていても, その交わりの上に余分な素因子は現れない.

このことは純正定理 \ref{thm:cm-iff-unmixed} ともきれいに対応している. 
実際, 周囲の滑らかな空間
$k[x,y,z]_{\widetilde{\mathfrak{m}}}$ の中で方程式 $xy=0$ は 1 本の非零因子で切られているので,
その零点集合は純粋に余次元 1 でなければならない. したがって現れる既約成分は 2 枚の平面
$V(x),V(y)$ だけであり, それらの交わりである $V(x,y)$ が余次元 2 の余分な成分として
現れることはない. これは局所的には 系\ref{cor:purity-local} の $r=1$ の場合にあたる. 
図\ref{fig:cm-example} で赤く描いた交線が「見えてはいるが成分ではない」のは, まさにこの純性の現れである. 

\vspace{2\baselineskip}

\begin{figure}[H]
    \centering
    \begin{tikzpicture}[
        line join=round,
        x={(1.18cm,0cm)},
        y={(-0.74cm,0.48cm)},
        z={(0cm,1.15cm)},
        info/.style={
            font=\small,
            fill=white,
            fill opacity=0.9,
            text opacity=1,
            rounded corners=2pt,
            inner sep=2pt
        },
        pointlabel/.style={
            font=\small,
            fill=white,
            fill opacity=0.9,
            text opacity=1,
            inner sep=1pt
        },
        callout/.style={
            ->,
            line width=0.9pt
        }
    ]
        \def\xmin{-2.2}
        \def\xmax{2.2}
        \def\ymin{-1.9}
        \def\ymax{1.9}
        \def\zmax{3.3}
        \def\zs{1.55}

        % ambient affine 3-space
        \draw[gray!55, dashed] (\xmin,\ymin,0) -- (\xmax,\ymin,0) -- (\xmax,\ymax,0) -- (\xmin,\ymax,0) -- cycle;
        \draw[gray!55, dashed] (\xmin,\ymin,\zmax) -- (\xmax,\ymin,\zmax) -- (\xmax,\ymax,\zmax) -- (\xmin,\ymax,\zmax) -- cycle;
        \draw[gray!55, dashed] (\xmin,\ymin,0) -- (\xmin,\ymin,\zmax);
        \draw[gray!55, dashed] (\xmax,\ymin,0) -- (\xmax,\ymin,\zmax);
        \draw[gray!55, dashed] (\xmax,\ymax,0) -- (\xmax,\ymax,\zmax);
        \draw[gray!55, dashed] (\xmin,\ymax,0) -- (\xmin,\ymax,\zmax);
        % horizontal slice showing the cross section xy=0
        \path[fill=gray!12, draw=gray!55, fill opacity=0.45]
            (\xmin,\ymin,\zs) -- (\xmax,\ymin,\zs) -- (\xmax,\ymax,\zs) -- (\xmin,\ymax,\zs) -- cycle;

        % coordinate planes
        \path[fill=blue!18, draw=blue!65!black, fill opacity=0.55]
            (0,\ymin,0) -- (0,\ymax,0) -- (0,\ymax,\zmax) -- (0,\ymin,\zmax) -- cycle;
        \path[fill=green!18, draw=green!50!black, fill opacity=0.55]
            (\xmin,0,0) -- (\xmax,0,0) -- (\xmax,0,\zmax) -- (\xmin,0,\zmax) -- cycle;
        \draw[green!18, line width=1.15pt] (0,\ymin,\zs) -- (1.17,0,\zs); 
        \draw[green!50!black, line width=1.15pt] (1.21,0,\zs) -- (\xmax,0,\zs);
        % intersection line
        \draw[red!70!black, line width=1.6pt] (0,0,0) -- (0,0,\zmax);

        % cross section on z=c
        \draw[blue!50!black, line width=1.2pt] (0,\ymin,\zmax) -- (0,\ymax,\zmax);
        \draw[blue!50!black, line width=1.2pt] (0,\ymin,0) -- (0,\ymin,\zmax);
        \draw[green!50!black, line width=1.2pt] (\xmin,0,\zs) -- (0,0,\zs);
        \draw[green!50!black, line width=1.2pt] (\xmin,0,0) -- (\xmin,0,\zmax);
        \draw[red!70!black, line width=1.2pt] (0,0,0) -- (0,0,\zmax); 
        \draw[green!18, line width=1.15pt] (0,0,\zs) -- (1.17,0,\zs); 
        \draw[green!50!black, line width=1.15pt] (1.21,0,\zs) -- (\xmax,0,\zs);

        % generic points
        \fill (0,1.20,2.30) circle (2.2pt);
        \fill (1.45,0,2.05) circle (2.2pt);
        \fill[red!70!black] (0,0,1.55) circle (2.5pt);

        % specialization arrows
        \draw[dashed, ->] (0.02,1.08,2.22) .. controls (0.02,0.72,2.00) and (0.01,0.35,1.80) .. (0.01,0.04,1.60);
        \draw[dashed, ->] (1.32,0.02,1.98) .. controls (0.95,0.02,1.86) and (0.48,0.02,1.72) .. (0.04,0.02,1.58);

        % callout targets
        \coordinate (vxTarget) at (0,1.35,2.75);
        \coordinate (vyTarget) at (1.85,0,2.75);
        \coordinate (vxyTarget) at (0,0,2.85);
        \coordinate (etaXTarget) at (0,1.20,2.30);
        \coordinate (etaYTarget) at (1.45,0,2.05);
        \coordinate (sliceTarget) at (1.10,0.40,\zs);
        \coordinate (nonassocTarget) at (0,0,1.55);

        % external labels with arrows
        \node[info, blue!65!black] (vxLabel) at (-3.95,0.98,3.02) {$V(x)$};
        \draw[callout, blue!65!black] (vxLabel.east) -- (vxTarget);

        \node[info] (etaXLabel) at (-3.65,0.20,2.35) {$\eta_{(x)}$};
        \draw[callout] (etaXLabel.east) -- (etaXTarget);

        \node[info, red!70!black] (vxyLabel) at (-0.10,2.55,4.00) {$V(x,y)$};
        \draw[callout, red!70!black] (vxyLabel.south) -- (vxyTarget);

        \node[info, green!50!black] (vyLabel) at (5.10,0.40,4.15) {$V(y)$};
        \draw[callout, green!50!black] (vyLabel.west) -- (vyTarget);

        \node[info, anchor=west] (etaYLabel) at (5.95,-0.70,1.70) {$\eta_{(y)}$};
        \draw[callout] (etaYLabel.west) .. controls (5.05,-0.55,1.76) and (3.40,-0.15,1.95) .. (etaYTarget);

        \node[info, red!70!black, align=left, anchor=west] (nonassocLabel) at (6.10,0.95,2.05) {$(x,y)$ は\\素因子ではない};
        \draw[callout, red!70!black] (nonassocLabel.west) -- (nonassocTarget);

        \node[info, align=left] (sliceLabel) at (5.40,-2.10,1.30) {断面 $z=c$ では\\$x=0$ と $y=0$ が十字に交差};
        \draw[callout, black!70] (sliceLabel.west) -- (sliceTarget);
    \end{tikzpicture}
    \caption{CM環 $R=k[x,y,z]/(xy)$ は 2 つの座標平面 $V(x),V(y)$ が交線 $V(x,y)$ に沿って交わる例である. 素因子は $(x),(y)$ のみであり, 交線に対応する $(x,y)$ は素因子ではなく, したがって非孤立素因子でもない.}
    \label{fig:cm-example}
\end{figure}

さらに scheme 的に見ると, $\Spec(R)$ の点は $R$ の素イデアルに対応している. 
図\ref{fig:cm-example} 中の $\eta_{(x)}$ は素イデアル $(x)$ に対応する点であり, その閉包が $V(x)$ なので
$V(x)$ の generic point を表している. ここで $\eta_{(x)}$ という記号は
「既約閉集合 $V(x)$ を代表する点」を表すための記号であって, 座標をもつ通常の点ではない.
実際, $\eta_{(x)}$ は閉点ではなく, その閉包がちょうど $V(x)$ 全体になる.
言い換えると, $V(x)$ のどの非空開集合にも $\eta_{(x)}$ は含まれるので, これが
$V(x)$ の最も generic な点である. 同様に $\eta_{(y)}$ は $V(y)$ の generic point である.
また赤い縦線は 2 つの平面の交線 $V(x,y)$ を表し, 赤い点は見た目には交線上の一点として
描いているが, 実際には素イデアル $(x,y)$ に対応する点, すなわち交線 $V(x,y)$ の generic point を
模式的に表している. この点が素因子ではないことが, 「\textbf{CM環では局所化して見ても非孤立素因子が現れない}」
という主張の具体例になっている. 

また図\ref{fig:cm-example} の灰色の切断面は各高さ $z=c$ の断面を表しており, そこで方程式 $xy=0$ は
2 本の直線 $x=0$ と $y=0$ の十字交差として見える.

\begin{rem}[非CMの場合との対比]
    ここまで見てきた「$\depth=\dim$」や「非孤立素因子が現れない」という性質が
    壊れると何が起きるかを見るために
    \[
        A:=k[x,y,z,w],\qquad I:=(x^2,xy)\subset A,\qquad S:=A/I
    \]
    を考える.  また
    \[
        \mathfrak{m}:=(x,y,z,w)/I \subset S,\qquad \mathfrak{p}:=(x,y)/I \subset S
    \]
    とおく. ここで $\mathfrak m$ は $S$ の極大イデアルであり,
    $\mathfrak p$ は $S$ の素イデアルである.
    なお $\mathfrak p$ は極大イデアルではない. 幾何学的には, $\mathfrak m$ は原点に対応し,
    $\mathfrak p$ は後で図に現れる 2 次元平面 $V(x,y)$ に対応する.
    このとき
    \[
        S_{\mathfrak{m}} = A_{(x,y,z,w)}/IA_{(x,y,z,w)}
    \]
    である. さらに
    \[
        \sqrt{I}=(x)\subset A
    \]
    である. したがって
    \[
        (S_{\mathfrak{m}})_{\mathrm{red}}
        \cong \bigl(A_{(x,y,z,w)}/IA_{(x,y,z,w)}\bigr)_{\mathrm{red}}
        \cong A_{(x,y,z,w)}/(x)
        \cong k[y,z,w]_{(y,z,w)}
    \]
    となる. 被約環にしても Krull 次元は変わらないので
    \[
        \dim S_{\mathfrak{m}}=\dim (S_{\mathfrak{m}})_{\mathrm{red}}=3
    \]
    である. 一方,
    \[
        S \cong \bigl(k[x,y]/(x^2,xy)\bigr)[z,w]
    \]
    なので, $z,w$ の像は $S$ 上の正則列であり, とくに $S_{\mathfrak{m}}$ 上の正則列である. 実際,
    \[
        S_{\mathfrak{m}}/(z,w)S_{\mathfrak{m}}
        \cong k[x,y]_{(x,y)}/(x^2,xy)
    \]
    であり, この商環では $\overline{x}\neq 0$ が $x$ と $y$ の両方で零化される.
    したがって正則元は存在せず,
    \[
        \depth\bigl(S_{\mathfrak{m}}/(z,w)S_{\mathfrak{m}}\bigr)=0
    \]
    となる. ゆえに
    \[
        \depth(S_{\mathfrak{m}})=2<3=\dim(S_{\mathfrak{m}})
    \]
    であり, 定義\ref{def:cm} の等式 $\depth=\dim$ が破れるので $S$ は CM環ではない.

    さらに
    \[
        I=(x)\cap(x^2,y)\qquad \text{in } A
    \]
    であるから, $\mathfrak{q}:=(x)/I \subset S$ が唯一の極小素因子であり,
    $\mathfrak{p}\subset S$ は非孤立素因子である.

    ここでは 命題\ref{prop:cm-associated-primes} の結論も壊れている.
    実際, $\mathfrak{p}\in \Ass_S(S)$ であるから局所化して
    \[
        \mathfrak{p}S_{\mathfrak{m}} \in \Ass_{S_{\mathfrak{m}}}(S_{\mathfrak{m}})
    \]
    となるが,
    \[
        S_{\mathfrak{m}}/\mathfrak{p}S_{\mathfrak{m}}
        \cong k[z,w]_{(z,w)}
    \]
    なので
    \[
        \dim\bigl(S_{\mathfrak{m}}/\mathfrak{p}S_{\mathfrak{m}}\bigr)=2
    \]
    である. しかし $\dim(S_{\mathfrak{m}})=3$ なので, もし $S_{\mathfrak{m}}$ が CM局所環なら
    命題\ref{prop:cm-associated-primes} により両者は一致するはずである. この不一致が,
    非CM性と非孤立素因子の出現を同時に示している.

    幾何学的には, $S$ の台集合は 4 次元空間の中の 3 次元超平面
    \[
        V(\mathfrak{q})=V(x)
    \]
    であり, その中の 2 次元平面
    \[
        V(\mathfrak{p})=V(x,y)
    \]
    に沿って nilpotent な厚みが乗っている.
    これに対して $R=k[x,y,z]/(xy)$ では, 2 枚の平面は交わっていても交線の generic point は素因子にならない.
    この違いが, CM環と非CM環の幾何の差をよく表している. 言い換えると, $S$ では 2 次元平面 $V(\mathfrak{p})=V(x,y)$ に対応する非孤立素因子が
    純正定理 \ref{thm:cm-iff-unmixed} の破れそのものとして現れているのに対し, $R$ では 系\ref{cor:purity-local} に沿って
    余分な成分が現れない.
\end{rem}

\vspace{2\baselineskip}

\begin{figure}[H]
    \centering
    \begin{tikzpicture}[
        line join=round,
        x={(1.18cm,0cm)},
        y={(-0.74cm,0.48cm)},
        z={(0cm,1.15cm)},
        info/.style={
            font=\small,
            fill=white,
            fill opacity=0.92,
            text opacity=1,
            rounded corners=2pt,
            inner sep=2pt
        },
        callout/.style={
            ->,
            line width=0.9pt
        }
    ]
        \def\ymaxA{2.55}
        \def\zmaxA{1.95}
        \def\wmaxA{3.15}
        \def\thickA{0.24}

        % support V(x) \cong A^3_{y,z,w}
        \path[fill=blue!14, fill opacity=0.24]
            (0,0,0) -- (\ymaxA,0,0) -- (\ymaxA,0,\wmaxA) -- (0,0,\wmaxA) -- cycle;
        \path[fill=blue!10, fill opacity=0.18]
            (\ymaxA,0,0) -- (\ymaxA,\zmaxA,0) -- (\ymaxA,\zmaxA,\wmaxA) -- (\ymaxA,0,\wmaxA) -- cycle;
        \path[fill=blue!8, fill opacity=0.14]
            (0,0,\wmaxA) -- (\ymaxA,0,\wmaxA) -- (\ymaxA,\zmaxA,\wmaxA) -- (0,\zmaxA,\wmaxA) -- cycle;

        % visible edges of A^3
        \draw[blue!65!black] (0,0,0) -- (\ymaxA,0,0) -- (\ymaxA,0,\wmaxA) -- (0,0,\wmaxA) -- cycle;
        \draw[blue!65!black] (\ymaxA,0,0) -- (\ymaxA,\zmaxA,0) -- (\ymaxA,\zmaxA,\wmaxA) -- (\ymaxA,0,\wmaxA);
        \draw[blue!65!black] (0,0,\wmaxA) -- (\ymaxA,0,\wmaxA) -- (\ymaxA,\zmaxA,\wmaxA) -- (0,\zmaxA,\wmaxA);

        % hidden edges of A^3
        \draw[blue!65!black, dashed] (0,0,0) -- (0,\zmaxA,0) -- (\ymaxA,\zmaxA,0);
        \draw[blue!65!black, dashed] (0,\zmaxA,0) -- (0,\zmaxA,\wmaxA);
        \draw[blue!65!black, dashed] (0,\zmaxA,\wmaxA) -- (0,0,\wmaxA);

        % infinitesimal thickening along V(x,y)
        \path[fill=red!20, fill opacity=0.18]
            (0,0,0) -- (\thickA,0,0) -- (\thickA,0,\wmaxA) -- (0,0,\wmaxA) -- cycle;
        \path[fill=red!16, fill opacity=0.14]
            (0,0,\wmaxA) -- (\thickA,0,\wmaxA) -- (\thickA,\zmaxA,\wmaxA) -- (0,\zmaxA,\wmaxA) -- cycle;
        \path[fill=red!22, draw=red!70!black, fill opacity=0.32]
            (0,0,0) -- (0,\zmaxA,0) -- (0,\zmaxA,\wmaxA) -- (0,0,\wmaxA) -- cycle;
        \path[fill=red!16, draw=red!70!black, fill opacity=0.14]
            (\thickA,0,0) -- (\thickA,\zmaxA,0) -- (\thickA,\zmaxA,\wmaxA) -- (\thickA,0,\wmaxA) -- cycle;
        \draw[red!70!black] (0,0,0) -- (\thickA,0,0);
        \draw[red!70!black] (0,0,\wmaxA) -- (\thickA,0,\wmaxA);
        \draw[red!70!black] (0,\zmaxA,\wmaxA) -- (\thickA,\zmaxA,\wmaxA);
        \draw[red!70!black, dashed] (0,\zmaxA,0) -- (\thickA,\zmaxA,0);

        % schematic points
        \fill (1.32,0.62,2.20) circle (2.2pt);
        \fill[red!70!black] (0.09,0.88,1.62) circle (2.5pt);

        % specialization arrow
        \draw[dashed, ->] (1.22,0.57,2.12) .. controls (0.96,0.56,2.00) and (0.45,0.78,1.82) .. (0.12,0.86,1.64);

        % axis labels inside V(x)
        \node[font=\scriptsize] at (\ymaxA+0.30,0.03,0) {$y$};
        \node[font=\scriptsize] at (-0.14,\zmaxA+0.18,0) {$z$};
        \node[font=\scriptsize] at (0.14,0.04,\wmaxA+0.25) {$w$};

        % callout targets
        \coordinate (spaceTarget) at (1.72,0.42,2.62);
        \coordinate (etaTarget) at (1.32,0.62,2.20);
        \coordinate (planeTarget) at (0,1.05,2.48);
        \coordinate (nonisolatedTarget) at (0.09,0.88,1.62);
        \coordinate (thickTarget) at (0.18,0.30,1.90);

        % external labels
        \node[info, blue!65!black, align=right] (spaceLabel) at (5.15,1.35,3.00) {$V(x)\cong \mathbb{A}^3_{y,z,w}$\\(唯一の既約成分)};
        \draw[callout, blue!65!black] (spaceLabel.west) -- (spaceTarget);

        \node[info] (etaLabel) at (-3.70,0.10,2.15) {$\eta_{(x)}$};
        \draw[callout] (etaLabel.east) -- (etaTarget);

        \node[info, red!70!black] (planeLabel) at (0.20,2.75,4.00) {$V(x,y)\cong \mathbb{A}^2_{z,w}$};
        \draw[callout, red!70!black] (planeLabel.south) -- (planeTarget);

        \node[info, red!70!black, align=left, anchor=west] (nonisolatedLabel) at (5.10,0.78,2.05) {$\mathfrak{p}$ は\\非孤立素因子};
        \draw[callout, red!70!black] (nonisolatedLabel.west) -- (nonisolatedTarget);

        \node[info, align=left, anchor=west] (thickLabel) at (5.00,-1.18,1.40) {赤い薄板が\\nilpotent な厚みを表す};
        \draw[callout] (thickLabel.west) -- (thickTarget);
    \end{tikzpicture}
    \caption{非CM環 $S=k[x,y,z,w]/(x^2,xy)$ の模式図. 台集合は 3 次元超平面 $V(x)\cong \mathbb{A}^3$ だけだが, その中の 2 次元平面 $V(x,y)=V(\mathfrak{p})$ に対応する非孤立素因子 $\mathfrak{p}=(x,y)/(x^2,xy)$ が現れ, そこに nilpotent な厚みが乗る.}
\end{figure}

図の赤い点は見た目の閉点ではなく, 2 次元平面 $V(x,y)=V(\mathfrak{p})$ の generic point を模式的に表している.

\section*{おわりに}

ここまで読んでいただきありがとうございます. 内容に関してもし何かお気づきの点があれば, ぜひご指摘いただけると嬉しいです. 今回は初等的な可換環論の話だけでCM環を定義して, その性質を見てきましたが, $\Ext$ などの derived な道具を使うと, CM環の性質がより自然に, そして簡明に見えるようになります. CM環より少し強く, 重要な概念として Gorenstein 環というものがありますが, 私の実力ではその説明をするのに derived な道具を使わざるを得ないので, そちらはまたの機会にでも...（先に $\Ext$ などを幾何的に特徴づける話をしたいと考えています.）

文献案内です. \cite{Matsumura} の17節の内容が今回の内容にあたります. 読むためにはある程度ホモロジー代数の知識が必要です. ほとんど読んでいませんが, この本に続く内容としては \cite{BrunsHerzog} が定番の本だと思います. scheme theoretic な内容は \cite{Liu} や \cite{Hartshorne} あたりが参考になると思います. 私は \cite{Liu} を読んでいますが, homological な言葉やcategorical な言葉を避けて書かれているので不満が少しあります. \cite{Liu} も \cite{Hartshorne} も CM環やGorenstein 環と関連した sheaf の面白い話はあまり出てこないので, そこは \cite{GoertzWedhorn} で補えます. これはhomological な言葉とcategorical な言葉を多用しているので、そういう言葉に慣れている人にはおすすめです. より代数幾何的な内容では, \cite{Badescu} が有名ですが, 私は読んでいません. 私の友人が面白いと言っていました. 

\vspace{1\baselineskip}

久しぶりに頑張って tikz で図を書いてみましたが, まだまだですね...ネタを吟味していたら時間がなくなり, 妥協した部分がかなりあります. 次回はもう少し計画的に書いてみたいと思います. ネタの候補はいくつか用意していますが, リクエストも歓迎します. （書けるかどうかは不明）
 

最後に, \underline{\Large{伊藤充洋}} \normalsize{および} \underline{\Large{新井秀斗} }\normalsize は東大に引きこもって数学をしてないで,  湧源のイベントやその他イベントに顔を出してください. 東大に引きこもられては私が寂しいです. puits なりでもいいので...

\begin{thebibliography}{99}

\bibitem{Matsumura}
松村 英之, \textit{可換環論}, 共立出版, 2000.

\bibitem{BrunsHerzog}
Winfried Bruns and J{\"u}rgen Herzog, \textit{Cohen-Macaulay Rings}, 2nd ed., Cambridge Studies in Advanced Mathematics 39, Cambridge University Press, 2009.

\bibitem{Liu}
Qing Liu, \textit{Algebraic Geometry and Arithmetic Curves}, Oxford University Press, 2002.

\bibitem{Hartshorne}
Robin Hartshorne, \textit{Algebraic Geometry}, Graduate Texts in Mathematics 52, Springer, 1977.

\bibitem{GoertzWedhorn}
Ulrich G{\"o}rtz and Torsten Wedhorn, \textit{Algebraic Geometry II: Cohomology of Schemes}, Springer Spektrum, Wiesbaden, 2023.

\bibitem{Badescu}
Lucian B{\u a}descu, \textit{Algebraic Surfaces}, Universitext, Springer, New York, 2001.

\end{thebibliography}

\end{document}
